% % % % % % % % % % % % % % % % % % % % % % % % % % % % % % % % % % % % % % % % % % % % % % % % % % % % % % % % % % % % % % % %
% FunctionalAnalysis 
% 1_04.tex
% 
% encoding: UTF-8 
% EOL: LF
%
% format: LaTeX
% indent: spaces (2)
% width: 127
% % % % % % % % % % % % % % % % % % % % % % % % % % % % % % % % % % % % % % % % % % % % % % % % % % % % % % % % % % % % % % % %
\textit{Let $B = \set[\big]{(z_1, z_2) \in \C^2}{\magnitude*{z_1} \leq \magnitude*{z_2}}$. Show that $B$ is balanced %
but that its interior is not. [compare with (e) of Theorem 1.13.]}
%
\begin{proof}
Let us equip $\C^2$ with the $1$-norm, that is, $\norm*[1]{(z_1, z_2)} = \magnitude*{z_1} + \magnitude*{z_2}$ for every pair %
$(z_1, z_2) \in \C^2$. Note that $B$ is nonempty, since the origin $(0, 0)$ lies in $B$. The proof consists of %
establishing that 
%
\begin{enumerate}
  \item $B$ is balanced,
  \item $B$ has nonempty interior,
  \item the interior $\interior{B}$ does not contain the origin. 
\end{enumerate}
%
This suffices, since every nonempty balanced subset contains the origin. Observe that
%
\begin{equation}
  \magnitude*{\alpha \cdot z_1}  
  = \magnitude*{\alpha}\cdot \magnitude*{z_1} 
  \leq  \magnitude*{\alpha} \cdot \magnitude*{z_2} 
  = \magnitude*{\alpha \cdot z_2} %
  \qquad \bigl((z_1,z_2) \in B, \alpha \in \C \bigr).
\end{equation}
%
In particular, this holds when $\magnitude*{\alpha} \leq 1$. Therefore, $B$ is balanced. This is (a). The triangle inequality %
gives
%
\begin{equation}
   \set[\big]{(z_1, z_2)}{ \norm*[1]{(1- z_1, 2- z_2)} < 1/3} \subset B. 
\end{equation}
%
Hence $\interior{B}\neq \emptyset$, which is (b). We now prove (c). To this end, suppose for the sake of contradiction that %
$B$ contains the origin-centered open ball
%
\begin{equation}
  U_r \Def \set{(z_1, z_2)}{\norm*[1]{(z_1, z_2)} < r}
\end{equation}
%
for some $r> 0$. The special case $(r/2, 0)\in U_r$ conflicts with the definition of $B$. This completes the proof. 
\end{proof}
%
\par\noindent{\textbf{Remark.}} This result implies that the assumption $0\in \interior{B}$ cannot be omitted from %
\citeFA{1.13 (e)}.
% END
% 

